\documentclass[11pt]{article}

    \usepackage[breakable]{tcolorbox}
    \usepackage{parskip} % Stop auto-indenting (to mimic markdown behaviour)
    

    % Basic figure setup, for now with no caption control since it's done
    % automatically by Pandoc (which extracts ![](path) syntax from Markdown).
    \usepackage{graphicx}
    % Keep aspect ratio if custom image width or height is specified
    \setkeys{Gin}{keepaspectratio}
    % Maintain compatibility with old templates. Remove in nbconvert 6.0
    \let\Oldincludegraphics\includegraphics
    % Ensure that by default, figures have no caption (until we provide a
    % proper Figure object with a Caption API and a way to capture that
    % in the conversion process - todo).
    \usepackage{caption}
    \DeclareCaptionFormat{nocaption}{}
    \captionsetup{format=nocaption,aboveskip=0pt,belowskip=0pt}

    \usepackage{float}
    \floatplacement{figure}{H} % forces figures to be placed at the correct location
    \usepackage{xcolor} % Allow colors to be defined
    \usepackage{enumerate} % Needed for markdown enumerations to work
    \usepackage{geometry} % Used to adjust the document margins
    \usepackage{amsmath} % Equations
    \usepackage{amssymb} % Equations
    \usepackage{textcomp} % defines textquotesingle
    \usepackage{amsmath}
    % Hack from http://tex.stackexchange.com/a/47451/13684:
    \AtBeginDocument{%
        \def\PYZsq{\textquotesingle}% Upright quotes in Pygmentized code
    }
    \usepackage{upquote} % Upright quotes for verbatim code
    \usepackage{eurosym} % defines \euro

    \usepackage{iftex}
    \ifPDFTeX
        \usepackage[T1]{fontenc}
        \IfFileExists{alphabeta.sty}{
              \usepackage{alphabeta}
          }{
              \usepackage[mathletters]{ucs}
              \usepackage[utf8x]{inputenc}
          }
    \else
        \usepackage{fontspec}
        \usepackage{unicode-math}
    \fi

    \usepackage{fancyvrb} % verbatim replacement that allows latex
    \usepackage{grffile} % extends the file name processing of package graphics
                         % to support a larger range
    \makeatletter % fix for old versions of grffile with XeLaTeX
    \@ifpackagelater{grffile}{2019/11/01}
    {
      % Do nothing on new versions
    }
    {
      \def\Gread@@xetex#1{%
        \IfFileExists{"\Gin@base".bb}%
        {\Gread@eps{\Gin@base.bb}}%
        {\Gread@@xetex@aux#1}%
      }
    }
    \makeatother
    \usepackage[Export]{adjustbox} % Used to constrain images to a maximum size
    \adjustboxset{max size={0.9\linewidth}{0.9\paperheight}}

    % The hyperref package gives us a pdf with properly built
    % internal navigation ('pdf bookmarks' for the table of contents,
    % internal cross-reference links, web links for URLs, etc.)
    \usepackage{hyperref}
    % The default LaTeX title has an obnoxious amount of whitespace. By default,
    % titling removes some of it. It also provides customization options.
    \usepackage{titling}
    \usepackage{longtable} % longtable support required by pandoc >1.10
    \usepackage{booktabs}  % table support for pandoc > 1.12.2
    \usepackage{array}     % table support for pandoc >= 2.11.3
    \usepackage{calc}      % table minipage width calculation for pandoc >= 2.11.1
    \usepackage[inline]{enumitem} % IRkernel/repr support (it uses the enumerate* environment)
    \usepackage[normalem]{ulem} % ulem is needed to support strikethroughs (\sout)
                                % normalem makes italics be italics, not underlines
    \usepackage{soul}      % strikethrough (\st) support for pandoc >= 3.0.0
    \usepackage{mathrsfs}
    

    
    % Colors for the hyperref package
    \definecolor{urlcolor}{rgb}{0,.145,.698}
    \definecolor{linkcolor}{rgb}{.71,0.21,0.01}
    \definecolor{citecolor}{rgb}{.12,.54,.11}

    % ANSI colors
    \definecolor{ansi-black}{HTML}{3E424D}
    \definecolor{ansi-black-intense}{HTML}{282C36}
    \definecolor{ansi-red}{HTML}{E75C58}
    \definecolor{ansi-red-intense}{HTML}{B22B31}
    \definecolor{ansi-green}{HTML}{00A250}
    \definecolor{ansi-green-intense}{HTML}{007427}
    \definecolor{ansi-yellow}{HTML}{DDB62B}
    \definecolor{ansi-yellow-intense}{HTML}{B27D12}
    \definecolor{ansi-blue}{HTML}{208FFB}
    \definecolor{ansi-blue-intense}{HTML}{0065CA}
    \definecolor{ansi-magenta}{HTML}{D160C4}
    \definecolor{ansi-magenta-intense}{HTML}{A03196}
    \definecolor{ansi-cyan}{HTML}{60C6C8}
    \definecolor{ansi-cyan-intense}{HTML}{258F8F}
    \definecolor{ansi-white}{HTML}{C5C1B4}
    \definecolor{ansi-white-intense}{HTML}{A1A6B2}
    \definecolor{ansi-default-inverse-fg}{HTML}{FFFFFF}
    \definecolor{ansi-default-inverse-bg}{HTML}{000000}

    % common color for the border for error outputs.
    \definecolor{outerrorbackground}{HTML}{FFDFDF}

    % commands and environments needed by pandoc snippets
    % extracted from the output of `pandoc -s`
    \providecommand{\tightlist}{%
      \setlength{\itemsep}{0pt}\setlength{\parskip}{0pt}}
    \DefineVerbatimEnvironment{Highlighting}{Verbatim}{commandchars=\\\{\}}
    % Add ',fontsize=\small' for more characters per line
    \newenvironment{Shaded}{}{}
    \newcommand{\KeywordTok}[1]{\textcolor[rgb]{0.00,0.44,0.13}{\textbf{{#1}}}}
    \newcommand{\DataTypeTok}[1]{\textcolor[rgb]{0.56,0.13,0.00}{{#1}}}
    \newcommand{\DecValTok}[1]{\textcolor[rgb]{0.25,0.63,0.44}{{#1}}}
    \newcommand{\BaseNTok}[1]{\textcolor[rgb]{0.25,0.63,0.44}{{#1}}}
    \newcommand{\FloatTok}[1]{\textcolor[rgb]{0.25,0.63,0.44}{{#1}}}
    \newcommand{\CharTok}[1]{\textcolor[rgb]{0.25,0.44,0.63}{{#1}}}
    \newcommand{\StringTok}[1]{\textcolor[rgb]{0.25,0.44,0.63}{{#1}}}
    \newcommand{\CommentTok}[1]{\textcolor[rgb]{0.38,0.63,0.69}{\textit{{#1}}}}
    \newcommand{\OtherTok}[1]{\textcolor[rgb]{0.00,0.44,0.13}{{#1}}}
    \newcommand{\AlertTok}[1]{\textcolor[rgb]{1.00,0.00,0.00}{\textbf{{#1}}}}
    \newcommand{\FunctionTok}[1]{\textcolor[rgb]{0.02,0.16,0.49}{{#1}}}
    \newcommand{\RegionMarkerTok}[1]{{#1}}
    \newcommand{\ErrorTok}[1]{\textcolor[rgb]{1.00,0.00,0.00}{\textbf{{#1}}}}
    \newcommand{\NormalTok}[1]{{#1}}

    % Additional commands for more recent versions of Pandoc
    \newcommand{\ConstantTok}[1]{\textcolor[rgb]{0.53,0.00,0.00}{{#1}}}
    \newcommand{\SpecialCharTok}[1]{\textcolor[rgb]{0.25,0.44,0.63}{{#1}}}
    \newcommand{\VerbatimStringTok}[1]{\textcolor[rgb]{0.25,0.44,0.63}{{#1}}}
    \newcommand{\SpecialStringTok}[1]{\textcolor[rgb]{0.73,0.40,0.53}{{#1}}}
    \newcommand{\ImportTok}[1]{{#1}}
    \newcommand{\DocumentationTok}[1]{\textcolor[rgb]{0.73,0.13,0.13}{\textit{{#1}}}}
    \newcommand{\AnnotationTok}[1]{\textcolor[rgb]{0.38,0.63,0.69}{\textbf{\textit{{#1}}}}}
    \newcommand{\CommentVarTok}[1]{\textcolor[rgb]{0.38,0.63,0.69}{\textbf{\textit{{#1}}}}}
    \newcommand{\VariableTok}[1]{\textcolor[rgb]{0.10,0.09,0.49}{{#1}}}
    \newcommand{\ControlFlowTok}[1]{\textcolor[rgb]{0.00,0.44,0.13}{\textbf{{#1}}}}
    \newcommand{\OperatorTok}[1]{\textcolor[rgb]{0.40,0.40,0.40}{{#1}}}
    \newcommand{\BuiltInTok}[1]{{#1}}
    \newcommand{\ExtensionTok}[1]{{#1}}
    \newcommand{\PreprocessorTok}[1]{\textcolor[rgb]{0.74,0.48,0.00}{{#1}}}
    \newcommand{\AttributeTok}[1]{\textcolor[rgb]{0.49,0.56,0.16}{{#1}}}
    \newcommand{\InformationTok}[1]{\textcolor[rgb]{0.38,0.63,0.69}{\textbf{\textit{{#1}}}}}
    \newcommand{\WarningTok}[1]{\textcolor[rgb]{0.38,0.63,0.69}{\textbf{\textit{{#1}}}}}
    \makeatletter
    \newsavebox\pandoc@box
    \newcommand*\pandocbounded[1]{%
      \sbox\pandoc@box{#1}%
      % scaling factors for width and height
      \Gscale@div\@tempa\textheight{\dimexpr\ht\pandoc@box+\dp\pandoc@box\relax}%
      \Gscale@div\@tempb\linewidth{\wd\pandoc@box}%
      % select the smaller of both
      \ifdim\@tempb\p@<\@tempa\p@
        \let\@tempa\@tempb
      \fi
      % scaling accordingly (\@tempa < 1)
      \ifdim\@tempa\p@<\p@
        \scalebox{\@tempa}{\usebox\pandoc@box}%
      % scaling not needed, use as it is
      \else
        \usebox{\pandoc@box}%
      \fi
    }
    \makeatother

    % Define a nice break command that doesn't care if a line doesn't already
    % exist.
    \def\br{\hspace*{\fill} \\* }
    % Math Jax compatibility definitions
    \def\gt{>}
    \def\lt{<}
    \let\Oldtex\TeX
    \let\Oldlatex\LaTeX
    \renewcommand{\TeX}{\textrm{\Oldtex}}
    \renewcommand{\LaTeX}{\textrm{\Oldlatex}}
    % Document parameters
    % Document title
    \title{Hughes\_CNET5050\_Assignment\_1\_Code}
    \author{by Ryan Hughes\\ Northeastern University\\ Master's in Complex Network Analysis}
    \date{}
    
    
    
    
    
    
% Pygments definitions
\makeatletter
\def\PY@reset{\let\PY@it=\relax \let\PY@bf=\relax%
    \let\PY@ul=\relax \let\PY@tc=\relax%
    \let\PY@bc=\relax \let\PY@ff=\relax}
\def\PY@tok#1{\csname PY@tok@#1\endcsname}
\def\PY@toks#1+{\ifx\relax#1\empty\else%
    \PY@tok{#1}\expandafter\PY@toks\fi}
\def\PY@do#1{\PY@bc{\PY@tc{\PY@ul{%
    \PY@it{\PY@bf{\PY@ff{#1}}}}}}}
\def\PY#1#2{\PY@reset\PY@toks#1+\relax+\PY@do{#2}}

\@namedef{PY@tok@w}{\def\PY@tc##1{\textcolor[rgb]{0.73,0.73,0.73}{##1}}}
\@namedef{PY@tok@c}{\let\PY@it=\textit\def\PY@tc##1{\textcolor[rgb]{0.24,0.48,0.48}{##1}}}
\@namedef{PY@tok@cp}{\def\PY@tc##1{\textcolor[rgb]{0.61,0.40,0.00}{##1}}}
\@namedef{PY@tok@k}{\let\PY@bf=\textbf\def\PY@tc##1{\textcolor[rgb]{0.00,0.50,0.00}{##1}}}
\@namedef{PY@tok@kp}{\def\PY@tc##1{\textcolor[rgb]{0.00,0.50,0.00}{##1}}}
\@namedef{PY@tok@kt}{\def\PY@tc##1{\textcolor[rgb]{0.69,0.00,0.25}{##1}}}
\@namedef{PY@tok@o}{\def\PY@tc##1{\textcolor[rgb]{0.40,0.40,0.40}{##1}}}
\@namedef{PY@tok@ow}{\let\PY@bf=\textbf\def\PY@tc##1{\textcolor[rgb]{0.67,0.13,1.00}{##1}}}
\@namedef{PY@tok@nb}{\def\PY@tc##1{\textcolor[rgb]{0.00,0.50,0.00}{##1}}}
\@namedef{PY@tok@nf}{\def\PY@tc##1{\textcolor[rgb]{0.00,0.00,1.00}{##1}}}
\@namedef{PY@tok@nc}{\let\PY@bf=\textbf\def\PY@tc##1{\textcolor[rgb]{0.00,0.00,1.00}{##1}}}
\@namedef{PY@tok@nn}{\let\PY@bf=\textbf\def\PY@tc##1{\textcolor[rgb]{0.00,0.00,1.00}{##1}}}
\@namedef{PY@tok@ne}{\let\PY@bf=\textbf\def\PY@tc##1{\textcolor[rgb]{0.80,0.25,0.22}{##1}}}
\@namedef{PY@tok@nv}{\def\PY@tc##1{\textcolor[rgb]{0.10,0.09,0.49}{##1}}}
\@namedef{PY@tok@no}{\def\PY@tc##1{\textcolor[rgb]{0.53,0.00,0.00}{##1}}}
\@namedef{PY@tok@nl}{\def\PY@tc##1{\textcolor[rgb]{0.46,0.46,0.00}{##1}}}
\@namedef{PY@tok@ni}{\let\PY@bf=\textbf\def\PY@tc##1{\textcolor[rgb]{0.44,0.44,0.44}{##1}}}
\@namedef{PY@tok@na}{\def\PY@tc##1{\textcolor[rgb]{0.41,0.47,0.13}{##1}}}
\@namedef{PY@tok@nt}{\let\PY@bf=\textbf\def\PY@tc##1{\textcolor[rgb]{0.00,0.50,0.00}{##1}}}
\@namedef{PY@tok@nd}{\def\PY@tc##1{\textcolor[rgb]{0.67,0.13,1.00}{##1}}}
\@namedef{PY@tok@s}{\def\PY@tc##1{\textcolor[rgb]{0.73,0.13,0.13}{##1}}}
\@namedef{PY@tok@sd}{\let\PY@it=\textit\def\PY@tc##1{\textcolor[rgb]{0.73,0.13,0.13}{##1}}}
\@namedef{PY@tok@si}{\let\PY@bf=\textbf\def\PY@tc##1{\textcolor[rgb]{0.64,0.35,0.47}{##1}}}
\@namedef{PY@tok@se}{\let\PY@bf=\textbf\def\PY@tc##1{\textcolor[rgb]{0.67,0.36,0.12}{##1}}}
\@namedef{PY@tok@sr}{\def\PY@tc##1{\textcolor[rgb]{0.64,0.35,0.47}{##1}}}
\@namedef{PY@tok@ss}{\def\PY@tc##1{\textcolor[rgb]{0.10,0.09,0.49}{##1}}}
\@namedef{PY@tok@sx}{\def\PY@tc##1{\textcolor[rgb]{0.00,0.50,0.00}{##1}}}
\@namedef{PY@tok@m}{\def\PY@tc##1{\textcolor[rgb]{0.40,0.40,0.40}{##1}}}
\@namedef{PY@tok@gh}{\let\PY@bf=\textbf\def\PY@tc##1{\textcolor[rgb]{0.00,0.00,0.50}{##1}}}
\@namedef{PY@tok@gu}{\let\PY@bf=\textbf\def\PY@tc##1{\textcolor[rgb]{0.50,0.00,0.50}{##1}}}
\@namedef{PY@tok@gd}{\def\PY@tc##1{\textcolor[rgb]{0.63,0.00,0.00}{##1}}}
\@namedef{PY@tok@gi}{\def\PY@tc##1{\textcolor[rgb]{0.00,0.52,0.00}{##1}}}
\@namedef{PY@tok@gr}{\def\PY@tc##1{\textcolor[rgb]{0.89,0.00,0.00}{##1}}}
\@namedef{PY@tok@ge}{\let\PY@it=\textit}
\@namedef{PY@tok@gs}{\let\PY@bf=\textbf}
\@namedef{PY@tok@ges}{\let\PY@bf=\textbf\let\PY@it=\textit}
\@namedef{PY@tok@gp}{\let\PY@bf=\textbf\def\PY@tc##1{\textcolor[rgb]{0.00,0.00,0.50}{##1}}}
\@namedef{PY@tok@go}{\def\PY@tc##1{\textcolor[rgb]{0.44,0.44,0.44}{##1}}}
\@namedef{PY@tok@gt}{\def\PY@tc##1{\textcolor[rgb]{0.00,0.27,0.87}{##1}}}
\@namedef{PY@tok@err}{\def\PY@bc##1{{\setlength{\fboxsep}{\string -\fboxrule}\fcolorbox[rgb]{1.00,0.00,0.00}{1,1,1}{\strut ##1}}}}
\@namedef{PY@tok@kc}{\let\PY@bf=\textbf\def\PY@tc##1{\textcolor[rgb]{0.00,0.50,0.00}{##1}}}
\@namedef{PY@tok@kd}{\let\PY@bf=\textbf\def\PY@tc##1{\textcolor[rgb]{0.00,0.50,0.00}{##1}}}
\@namedef{PY@tok@kn}{\let\PY@bf=\textbf\def\PY@tc##1{\textcolor[rgb]{0.00,0.50,0.00}{##1}}}
\@namedef{PY@tok@kr}{\let\PY@bf=\textbf\def\PY@tc##1{\textcolor[rgb]{0.00,0.50,0.00}{##1}}}
\@namedef{PY@tok@bp}{\def\PY@tc##1{\textcolor[rgb]{0.00,0.50,0.00}{##1}}}
\@namedef{PY@tok@fm}{\def\PY@tc##1{\textcolor[rgb]{0.00,0.00,1.00}{##1}}}
\@namedef{PY@tok@vc}{\def\PY@tc##1{\textcolor[rgb]{0.10,0.09,0.49}{##1}}}
\@namedef{PY@tok@vg}{\def\PY@tc##1{\textcolor[rgb]{0.10,0.09,0.49}{##1}}}
\@namedef{PY@tok@vi}{\def\PY@tc##1{\textcolor[rgb]{0.10,0.09,0.49}{##1}}}
\@namedef{PY@tok@vm}{\def\PY@tc##1{\textcolor[rgb]{0.10,0.09,0.49}{##1}}}
\@namedef{PY@tok@sa}{\def\PY@tc##1{\textcolor[rgb]{0.73,0.13,0.13}{##1}}}
\@namedef{PY@tok@sb}{\def\PY@tc##1{\textcolor[rgb]{0.73,0.13,0.13}{##1}}}
\@namedef{PY@tok@sc}{\def\PY@tc##1{\textcolor[rgb]{0.73,0.13,0.13}{##1}}}
\@namedef{PY@tok@dl}{\def\PY@tc##1{\textcolor[rgb]{0.73,0.13,0.13}{##1}}}
\@namedef{PY@tok@s2}{\def\PY@tc##1{\textcolor[rgb]{0.73,0.13,0.13}{##1}}}
\@namedef{PY@tok@sh}{\def\PY@tc##1{\textcolor[rgb]{0.73,0.13,0.13}{##1}}}
\@namedef{PY@tok@s1}{\def\PY@tc##1{\textcolor[rgb]{0.73,0.13,0.13}{##1}}}
\@namedef{PY@tok@mb}{\def\PY@tc##1{\textcolor[rgb]{0.40,0.40,0.40}{##1}}}
\@namedef{PY@tok@mf}{\def\PY@tc##1{\textcolor[rgb]{0.40,0.40,0.40}{##1}}}
\@namedef{PY@tok@mh}{\def\PY@tc##1{\textcolor[rgb]{0.40,0.40,0.40}{##1}}}
\@namedef{PY@tok@mi}{\def\PY@tc##1{\textcolor[rgb]{0.40,0.40,0.40}{##1}}}
\@namedef{PY@tok@il}{\def\PY@tc##1{\textcolor[rgb]{0.40,0.40,0.40}{##1}}}
\@namedef{PY@tok@mo}{\def\PY@tc##1{\textcolor[rgb]{0.40,0.40,0.40}{##1}}}
\@namedef{PY@tok@ch}{\let\PY@it=\textit\def\PY@tc##1{\textcolor[rgb]{0.24,0.48,0.48}{##1}}}
\@namedef{PY@tok@cm}{\let\PY@it=\textit\def\PY@tc##1{\textcolor[rgb]{0.24,0.48,0.48}{##1}}}
\@namedef{PY@tok@cpf}{\let\PY@it=\textit\def\PY@tc##1{\textcolor[rgb]{0.24,0.48,0.48}{##1}}}
\@namedef{PY@tok@c1}{\let\PY@it=\textit\def\PY@tc##1{\textcolor[rgb]{0.24,0.48,0.48}{##1}}}
\@namedef{PY@tok@cs}{\let\PY@it=\textit\def\PY@tc##1{\textcolor[rgb]{0.24,0.48,0.48}{##1}}}

\def\PYZbs{\char`\\}
\def\PYZus{\char`\_}
\def\PYZob{\char`\{}
\def\PYZcb{\char`\}}
\def\PYZca{\char`\^}
\def\PYZam{\char`\&}
\def\PYZlt{\char`\<}
\def\PYZgt{\char`\>}
\def\PYZsh{\char`\#}
\def\PYZpc{\char`\%}
\def\PYZdl{\char`\$}
\def\PYZhy{\char`\-}
\def\PYZsq{\char`\'}
\def\PYZdq{\char`\"}
\def\PYZti{\char`\~}
% for compatibility with earlier versions
\def\PYZat{@}
\def\PYZlb{[}
\def\PYZrb{]}
\makeatother


    % For linebreaks inside Verbatim environment from package fancyvrb.
    \makeatletter
        \newbox\Wrappedcontinuationbox
        \newbox\Wrappedvisiblespacebox
        \newcommand*\Wrappedvisiblespace {\textcolor{red}{\textvisiblespace}}
        \newcommand*\Wrappedcontinuationsymbol {\textcolor{red}{\llap{\tiny$\m@th\hookrightarrow$}}}
        \newcommand*\Wrappedcontinuationindent {3ex }
        \newcommand*\Wrappedafterbreak {\kern\Wrappedcontinuationindent\copy\Wrappedcontinuationbox}
        % Take advantage of the already applied Pygments mark-up to insert
        % potential linebreaks for TeX processing.
        %        {, <, #, %, $, ' and ": go to next line.
        %        _, }, ^, &, >, - and ~: stay at end of broken line.
        % Use of \textquotesingle for straight quote.
        \newcommand*\Wrappedbreaksatspecials {%
            \def\PYGZus{\discretionary{\char`\_}{\Wrappedafterbreak}{\char`\_}}%
            \def\PYGZob{\discretionary{}{\Wrappedafterbreak\char`\{}{\char`\{}}%
            \def\PYGZcb{\discretionary{\char`\}}{\Wrappedafterbreak}{\char`\}}}%
            \def\PYGZca{\discretionary{\char`\^}{\Wrappedafterbreak}{\char`\^}}%
            \def\PYGZam{\discretionary{\char`\&}{\Wrappedafterbreak}{\char`\&}}%
            \def\PYGZlt{\discretionary{}{\Wrappedafterbreak\char`\<}{\char`\<}}%
            \def\PYGZgt{\discretionary{\char`\>}{\Wrappedafterbreak}{\char`\>}}%
            \def\PYGZsh{\discretionary{}{\Wrappedafterbreak\char`\#}{\char`\#}}%
            \def\PYGZpc{\discretionary{}{\Wrappedafterbreak\char`\%}{\char`\%}}%
            \def\PYGZdl{\discretionary{}{\Wrappedafterbreak\char`\$}{\char`\$}}%
            \def\PYGZhy{\discretionary{\char`\-}{\Wrappedafterbreak}{\char`\-}}%
            \def\PYGZsq{\discretionary{}{\Wrappedafterbreak\textquotesingle}{\textquotesingle}}%
            \def\PYGZdq{\discretionary{}{\Wrappedafterbreak\char`\"}{\char`\"}}%
            \def\PYGZti{\discretionary{\char`\~}{\Wrappedafterbreak}{\char`\~}}%
        }
        % Some characters . , ; ? ! / are not pygmentized.
        % This macro makes them "active" and they will insert potential linebreaks
        \newcommand*\Wrappedbreaksatpunct {%
            \lccode`\~`\.\lowercase{\def~}{\discretionary{\hbox{\char`\.}}{\Wrappedafterbreak}{\hbox{\char`\.}}}%
            \lccode`\~`\,\lowercase{\def~}{\discretionary{\hbox{\char`\,}}{\Wrappedafterbreak}{\hbox{\char`\,}}}%
            \lccode`\~`\;\lowercase{\def~}{\discretionary{\hbox{\char`\;}}{\Wrappedafterbreak}{\hbox{\char`\;}}}%
            \lccode`\~`\:\lowercase{\def~}{\discretionary{\hbox{\char`\:}}{\Wrappedafterbreak}{\hbox{\char`\:}}}%
            \lccode`\~`\?\lowercase{\def~}{\discretionary{\hbox{\char`\?}}{\Wrappedafterbreak}{\hbox{\char`\?}}}%
            \lccode`\~`\!\lowercase{\def~}{\discretionary{\hbox{\char`\!}}{\Wrappedafterbreak}{\hbox{\char`\!}}}%
            \lccode`\~`\/\lowercase{\def~}{\discretionary{\hbox{\char`\/}}{\Wrappedafterbreak}{\hbox{\char`\/}}}%
            \catcode`\.\active
            \catcode`\,\active
            \catcode`\;\active
            \catcode`\:\active
            \catcode`\?\active
            \catcode`\!\active
            \catcode`\/\active
            \lccode`\~`\~
        }
    \makeatother

    \let\OriginalVerbatim=\Verbatim
    \makeatletter
    \renewcommand{\Verbatim}[1][1]{%
        %\parskip\z@skip
        \sbox\Wrappedcontinuationbox {\Wrappedcontinuationsymbol}%
        \sbox\Wrappedvisiblespacebox {\FV@SetupFont\Wrappedvisiblespace}%
        \def\FancyVerbFormatLine ##1{\hsize\linewidth
            \vtop{\raggedright\hyphenpenalty\z@\exhyphenpenalty\z@
                \doublehyphendemerits\z@\finalhyphendemerits\z@
                \strut ##1\strut}%
        }%
        % If the linebreak is at a space, the latter will be displayed as visible
        % space at end of first line, and a continuation symbol starts next line.
        % Stretch/shrink are however usually zero for typewriter font.
        \def\FV@Space {%
            \nobreak\hskip\z@ plus\fontdimen3\font minus\fontdimen4\font
            \discretionary{\copy\Wrappedvisiblespacebox}{\Wrappedafterbreak}
            {\kern\fontdimen2\font}%
        }%

        % Allow breaks at special characters using \PYG... macros.
        \Wrappedbreaksatspecials
        % Breaks at punctuation characters . , ; ? ! and / need catcode=\active
        \OriginalVerbatim[#1,codes*=\Wrappedbreaksatpunct]%
    }
    \makeatother

    % Exact colors from NB
    \definecolor{incolor}{HTML}{303F9F}
    \definecolor{outcolor}{HTML}{D84315}
    \definecolor{cellborder}{HTML}{CFCFCF}
    \definecolor{cellbackground}{HTML}{F7F7F7}

    % prompt
    \makeatletter
    \newcommand{\boxspacing}{\kern\kvtcb@left@rule\kern\kvtcb@boxsep}
    \makeatother
    \newcommand{\prompt}[4]{
        {\ttfamily\llap{{\color{#2}[#3]:\hspace{3pt}#4}}\vspace{-\baselineskip}}
    }
    

    
    % Prevent overflowing lines due to hard-to-break entities
    \sloppy
    % Setup hyperref package
    \hypersetup{
      breaklinks=true,  % so long urls are correctly broken across lines
      colorlinks=true,
      urlcolor=urlcolor,
      linkcolor=linkcolor,
      citecolor=citecolor,
      }
    % Slightly bigger margins than the latex defaults
    
    \geometry{verbose,tmargin=1in,bmargin=1in,lmargin=1in,rmargin=1in}
    
    

\begin{document}
    
    \maketitle


    
    \paragraph{\newline Introduction.}\label{introduction.}


In Part 1 of this assignment, you will analyze selected properties of
random and real-world networks. Submit your written solutions as a PDF
prepared using LaTeX or Microsoft Word.

You must also submit all code used to complete the assignment, either as
separate files or as clearly formatted appendices in your write-up.

Ensure that your code is readable and appropriately commented. If you
use generative AI tools such as ChatGPT, clearly identify the tasks for
which they were used and include the relevant prompts and generated
outputs.

In Part 2, you will develop your ability to critically read and present
network science research papers.

\paragraph{\newline Part 1}\label{part-1}

\paragraph{I. Random networks {[}25
marks{]}}\label{i.-random-networks-25-marks}

\paragraph{(a) Consider the Erdos--Renyi random graph model G(n, p) with n =
1000 nodes and edge probability p = 0.01. Write down the probability
distribution P(L = l) for the total number of edges L, and explain each
term in your expression. Then calculate the expected number of edges,
⟨L⟩, for the specified values of n and
p.\newline \newline}\label{consider-the-erdosrenyi-random-graph-model-gn-p-with-n-1000-nodes-and-edge-probability-p-0.01.a-write-down-the-probability-distribution-pl-l-for-the-total-number-of-edges-l-and-explain-each-term-in-your-expression.-then-calculate-the-expected-number-of-edges-l-for-the-specified-values-of-n-and-p.}


Probability Distribution:


\[ P(L)=  {\frac{N(N-1)}{2}\choose L}p^L (1-p)^{\frac{N(N-1)}{2}-L} \]\newline


This equation reflects the probability distribution for some probability
\(p\) , that there are exactly `\(L\) = l' total number of edges in an
Erdos--Renyi random graph model.\newline


To explain further, there are 4 variables represented here:
\(P(L)\), the probability for \(L\) edges. \(L\), the number of edges. \(N\), the number of nodes.
And, \(p\), the probability of an edge existing between any pair of nodes.

These variables make up the following key terms of the equation: \newline
\(p^L\) reflects the probability of drawing an edge an arbitrary
amount of times equal to the value of \(L\)

\((1-p)^{\frac{N(N-1)}{2}-L}\) reflects the probability we don't draw an edge
 for all other pairs of nodes that one could place an edge.\newline

Or, in other words: the
probability we don't make an edge over the remaining ``N choose two
(\(\frac{N(N-1)}{2}\))'' pairs leftover after \(L\) pairs have been
processed. And,

\({\frac{N(N-1)}{2}\choose L}\) represents ``N choose 2, choose L'' or,
the picking of 2 unique nodes, an arbitrary amount of times: \(L\).

\paragraph{Calculating the expected number of edges for
G(1000,0.01):\newline}\label{calculating-the-expected-number-of-edges-for-g10000.01}



 $\newline$
  \[ \langle{L}\rangle = p\frac{N(N-1)}{2} \]

 \[ \langle{L}\rangle = 0.01\frac{1000(1000-1)}{2} \]

 \[ \langle{L}\rangle = 0.01*\frac{1000(999)}{2} \]

 \[ \langle{L}\rangle = 0.01*499500 \]

 \[ \langle{L}\rangle = 4995 \]

\paragraph{(b) Write down the degree distribution P(K = k) and explain each
term in your expression. Calculate the expected degree, \(\langle{k}\rangle\), for the
specified values of n and
p.\newline \newline}\label{write-down-the-degree-distribution-pk-k-and-explain-each-term-in-your-expression.calculate-the-expected-degree-k-for-the-specified-values-of-n-and-p.}

Degree Distribution:

\[ p_k=  {{(N-1)}\choose k}p^k (1-p)^{(N-1)-k} \] $\newline$

This equation reflects the probability distribution for some probability
\(p\) , that there are exactly `\(K\) = k' total degree in an
Erdos--Renyi random graph model.

To explain further, there are just 3 variables in this equation.
\(p_k\), the probability for \(k\) degree to occur for a given node. \(N\), the number of nodes.
And, \(p\), the probability of an edge existing between any pair of nodes.

These variables make up the following key terms of the equation: \newline
\(p^k\) reflects the probability of a node having a degree of an arbitrary
amount of times, equal to the value of \(k\)

\((1-p)^{(N-1)-k}\) in a similar fashion, reflects the probability we
don't have an edge (i.e., an edge contributing to the degree of a given node) for all other possibilities available.

\({{(N-1)}\choose k}\) represents ``N-1 choose k'' or, the picking of
one other node, an arbitrary amount of times equal to the value of \(k\).\newline



Calculating the expected average degree for G(1000,0.01):

 $\newline$
 \[ \langle{k}\rangle = p(N-1) \]

 \[  \langle{k}\rangle = 0.01(1000-1) \]

 \[  \langle{k}\rangle = 0.01*(999) \]

 \[  \langle{k}\rangle = 9.99 \]

\paragraph{(c) Write down the definition for average local clustering
coefficient. What is the average local clustering coefficient of G(n, p),
for the specified values of n and
p?\newline \newline}\label{c-write-down-the-definition-for-average-local-clustering-coefficient.-what-is-the-average-localclustering-coefficient-of-gn-p-for-the-specified-values-of-n-and-p}

\(C_i = p\), is the definition for the average local clustering
coefficient in Erdos-Renyi Random Graph models. Due to the fact, the
placing of any edge is simply \(p\), the odds of a node's neighbors
being connected (and thereby, there being clustering) is also \(p\).

As such, the average local clustering coefficient for \(G(1000,0.01)\)
is:

\[C_i = p\]

\[C_i = 0.01\]

\paragraph{(d) Write down the approximate expression for the average
shortest-path length \(\langle{d}\rangle\) between nodes in G(n, p) and calculate it for
the specified values of n and
p.\newline \newline}\label{write-down-the-approximate-expression-for-the-average-shortest-path-length-d-between-nodes-in-gn-p-and-calculate-it-for-the-specified-values-of-n-and-p.}

\(d \approx \frac{logN}{log\langle{k}\rangle}\), is the approximate
expression for the average shortest-path length, \(\langle{d}\rangle\),
between nodes in a graph \(G(N,p)\).\newline

Therefore, for \(G(N,p) = G(1000,0.01)\):

\[d \approx \frac{logN}{log\langle{k}\rangle}\]

\[d \approx \frac{log1000}{log9.99}\]

\[d \approx \frac{3}{0.9995655}\]

\[d \approx 3.001304\]

\paragraph{(e) Write a program in Python, or another programming language of
your choice, that takes n and p as inputs and returns one realization of
G(n,p).}\label{write-a-program-in-python-or-another-programming-language-of-your-choice-that-takes-nand-p-as-inputs-and-returns-one-realization-of-gn-p.}

    \begin{tcolorbox}[breakable, size=fbox, boxrule=1pt, pad at break*=1mm,colback=cellbackground, colframe=cellborder]
\prompt{In}{incolor}{16}{\boxspacing}
\begin{Verbatim}[commandchars=\\\{\}]
\PY{k+kn}{from}\PY{+w}{ }\PY{n+nn}{itertools}\PY{+w}{ }\PY{k+kn}{import} \PY{n}{combinations}
\PY{k+kn}{import}\PY{+w}{ }\PY{n+nn}{networkx}\PY{+w}{ }\PY{k}{as}\PY{+w}{ }\PY{n+nn}{nx}
\PY{k+kn}{import}\PY{+w}{ }\PY{n+nn}{numpy}\PY{+w}{ }\PY{k}{as}\PY{+w}{ }\PY{n+nn}{np}
\PY{k+kn}{import}\PY{+w}{ }\PY{n+nn}{scipy}\PY{n+nn}{.}\PY{n+nn}{stats}


\PY{k}{def}\PY{+w}{ }\PY{n+nf}{create\PYZus{}erdos\PYZus{}renyi\PYZus{}random\PYZus{}graph}\PY{p}{(}\PY{n}{N}\PY{p}{,} \PY{n}{p}\PY{p}{)}\PY{p}{:}
\PY{+w}{    }\PY{l+s+sd}{\PYZdq{}\PYZdq{}\PYZdq{}}
\PY{l+s+sd}{    Creates a Erdos\PYZhy{}Renyi Random Graph with \PYZsq{}N\PYZsq{} nodes and an edge probability, \PYZsq{}p\PYZsq{}}

\PY{l+s+sd}{    Parameters}
\PY{l+s+sd}{    \PYZhy{}\PYZhy{}\PYZhy{}\PYZhy{}\PYZhy{}\PYZhy{}\PYZhy{}\PYZhy{}\PYZhy{}\PYZhy{}}
\PY{l+s+sd}{    N : int}
\PY{l+s+sd}{        The number of nodes of the graph to be created.}

\PY{l+s+sd}{    p : float}
\PY{l+s+sd}{        The probability an edge is placed between any pair of nodes}

\PY{l+s+sd}{    Returns}
\PY{l+s+sd}{    \PYZhy{}\PYZhy{}\PYZhy{}\PYZhy{}\PYZhy{}\PYZhy{}\PYZhy{}}
\PY{l+s+sd}{    G : networkx.Graph}
\PY{l+s+sd}{        The resulting erdos\PYZhy{}renyi random graph,undirected, as supported by NetworkX.}

\PY{l+s+sd}{    \PYZdq{}\PYZdq{}\PYZdq{}}
    \PY{c+c1}{\PYZsh{}Initialize the Graph}
    \PY{k+kn}{from}\PY{+w}{ }\PY{n+nn}{itertools}\PY{+w}{ }\PY{k+kn}{import} \PY{n}{combinations}
    \PY{k+kn}{import}\PY{+w}{ }\PY{n+nn}{networkx}\PY{+w}{ }\PY{k}{as}\PY{+w}{ }\PY{n+nn}{nx}
    \PY{n}{G} \PY{o}{=} \PY{n}{nx}\PY{o}{.}\PY{n}{Graph}\PY{p}{(}\PY{p}{)}
    \PY{c+c1}{\PYZsh{}Initialize chance}
    \PY{n}{chance} \PY{o}{=} \PY{l+m+mi}{0}

    \PY{c+c1}{\PYZsh{}Add in all the nodes}
    \PY{k}{for} \PY{n}{i} \PY{o+ow}{in} \PY{n+nb}{range}\PY{p}{(}\PY{n}{N}\PY{p}{)}\PY{p}{:}
        \PY{c+c1}{\PYZsh{}Add a node}
        \PY{n}{G}\PY{o}{.}\PY{n}{add\PYZus{}node}\PY{p}{(}\PY{n}{i}\PY{p}{)}
    \PY{c+c1}{\PYZsh{}Add in all the edges}
    \PY{k}{for} \PY{n}{node1\PYZus{}id}\PY{p}{,}\PY{n}{node2\PYZus{}id} \PY{o+ow}{in} \PY{n}{combinations}\PY{p}{(}\PY{n}{G}\PY{o}{.}\PY{n}{nodes}\PY{p}{(}\PY{p}{)}\PY{p}{,} \PY{l+m+mi}{2}\PY{p}{)}\PY{p}{:}\PY{c+c1}{\PYZsh{}For every pair of nodes within G}
        \PY{n}{chance} \PY{o}{=} \PY{n}{np}\PY{o}{.}\PY{n}{random}\PY{o}{.}\PY{n}{rand}\PY{p}{(}\PY{p}{)}\PY{c+c1}{\PYZsh{}Draw a random chance 0\PYZhy{}100\PYZpc{}}
        \PY{k}{if} \PY{n}{chance} \PY{o}{\PYZlt{}}\PY{o}{=} \PY{n}{p}\PY{p}{:}\PY{c+c1}{\PYZsh{}If the chance is within the threshold \PYZsq{}p\PYZsq{}}
            \PY{c+c1}{\PYZsh{}Add an edge between that pair of nodes}
            \PY{n}{G}\PY{o}{.}\PY{n}{add\PYZus{}edge}\PY{p}{(}\PY{n}{node1\PYZus{}id}\PY{p}{,}\PY{n}{node2\PYZus{}id}\PY{p}{)}
    \PY{k}{return} \PY{n}{G}

\PY{n}{N} \PY{o}{=} \PY{l+m+mi}{100}
\PY{n}{p} \PY{o}{=} \PY{l+m+mf}{0.01}
\PY{n+nb}{print}\PY{p}{(}\PY{n}{nx}\PY{o}{.}\PY{n}{erdos\PYZus{}renyi\PYZus{}graph}\PY{p}{(}\PY{l+m+mi}{100}\PY{p}{,}\PY{l+m+mf}{0.01}\PY{p}{)}\PY{p}{)}
\PY{n+nb}{print}\PY{p}{(}\PY{n}{create\PYZus{}erdos\PYZus{}renyi\PYZus{}random\PYZus{}graph}\PY{p}{(}\PY{n}{N}\PY{p}{,}\PY{n}{p}\PY{p}{)}\PY{p}{)}
\end{Verbatim}
\end{tcolorbox}

    \begin{Verbatim}[commandchars=\\\{\}]
Graph with 100 nodes and 57 edges
Graph with 100 nodes and 62 edges
    \end{Verbatim}

    \paragraph{(i) Generate NG = 100 independent graph realizations. For
each realization, record the total number of edges. Plot the empirical
distribution of these 100 edge counts and overlay the theoretical
distribution derived in part (a). Clearly label the axes
and legend.}\label{i-generate-ng-100-independent-graph-realizations.-for-each-realization-record-thetotal-number-of-edges.-plot-the-empirical-distribution-of-these-100-edge-counts-andoverlay-the-theoretical-distribution-derived-in-part-a.-clearly-label-the-axes-andlegend.}

    \begin{tcolorbox}[breakable, size=fbox, boxrule=1pt, pad at break*=1mm,colback=cellbackground, colframe=cellborder]
\prompt{In}{incolor}{17}{\boxspacing}
\begin{Verbatim}[commandchars=\\\{\}]
\PY{k+kn}{import}\PY{+w}{ }\PY{n+nn}{networkx}\PY{+w}{ }\PY{k}{as}\PY{+w}{ }\PY{n+nn}{nx}
\PY{k+kn}{import}\PY{+w}{ }\PY{n+nn}{matplotlib}\PY{n+nn}{.}\PY{n+nn}{pyplot}\PY{+w}{ }\PY{k}{as}\PY{+w}{ }\PY{n+nn}{plt}
\PY{k+kn}{from}\PY{+w}{ }\PY{n+nn}{scipy}\PY{n+nn}{.}\PY{n+nn}{stats}\PY{+w}{ }\PY{k+kn}{import} \PY{n}{binom}
\PY{k+kn}{from}\PY{+w}{ }\PY{n+nn}{scipy}\PY{n+nn}{.}\PY{n+nn}{special}\PY{+w}{ }\PY{k+kn}{import} \PY{n}{comb}

\PY{n}{N}\PY{o}{=}\PY{l+m+mi}{1000}
\PY{n}{p}\PY{o}{=}\PY{l+m+mf}{0.01}
\PY{n}{num\PYZus{}edges\PYZus{}list}\PY{o}{=}\PY{p}{[}\PY{p}{]}
\PY{n}{my\PYZus{}num\PYZus{}edges\PYZus{}list}\PY{o}{=}\PY{p}{[}\PY{p}{]}
\PY{k}{for} \PY{n}{i} \PY{o+ow}{in} \PY{n+nb}{range}\PY{p}{(}\PY{l+m+mi}{100}\PY{p}{)}\PY{p}{:}
  \PY{n}{G}\PY{o}{=}\PY{n}{nx}\PY{o}{.}\PY{n}{erdos\PYZus{}renyi\PYZus{}graph}\PY{p}{(}\PY{n}{N}\PY{p}{,} \PY{n}{p}\PY{p}{)}
  \PY{n}{edges} \PY{o}{=} \PY{n+nb}{list}\PY{p}{(}\PY{n}{G}\PY{o}{.}\PY{n}{copy}\PY{p}{(}\PY{p}{)}\PY{o}{.}\PY{n}{edges}\PY{p}{(}\PY{p}{)}\PY{p}{)}
  \PY{n}{num\PYZus{}edges\PYZus{}list}\PY{o}{.}\PY{n}{append}\PY{p}{(}\PY{n+nb}{len}\PY{p}{(}\PY{n}{edges}\PY{p}{)}\PY{p}{)}
  \PY{n}{G2}\PY{o}{=}\PY{n}{create\PYZus{}erdos\PYZus{}renyi\PYZus{}random\PYZus{}graph}\PY{p}{(}\PY{n}{N}\PY{p}{,} \PY{n}{p}\PY{p}{)}
  \PY{n}{edges} \PY{o}{=} \PY{n+nb}{list}\PY{p}{(}\PY{n}{G2}\PY{o}{.}\PY{n}{copy}\PY{p}{(}\PY{p}{)}\PY{o}{.}\PY{n}{edges}\PY{p}{(}\PY{p}{)}\PY{p}{)}
  \PY{n}{my\PYZus{}num\PYZus{}edges\PYZus{}list}\PY{o}{.}\PY{n}{append}\PY{p}{(}\PY{n+nb}{len}\PY{p}{(}\PY{n}{edges}\PY{p}{)}\PY{p}{)}

\PY{n}{L} \PY{o}{=} \PY{n}{np}\PY{o}{.}\PY{n}{arange}\PY{p}{(}\PY{l+m+mi}{4995}\PY{o}{\PYZhy{}}\PY{l+m+mi}{200}\PY{p}{,}\PY{l+m+mi}{4995}\PY{o}{+}\PY{l+m+mi}{200}\PY{o}{+}\PY{l+m+mi}{1}\PY{p}{,}\PY{n}{step}\PY{o}{=}\PY{l+m+mi}{1}\PY{p}{,}\PY{n}{dtype} \PY{o}{=} \PY{p}{(}\PY{n+nb}{int}\PY{p}{)}\PY{p}{)} \PY{c+c1}{\PYZsh{}Picking a handful of values centered around the average I calculated earlier of: 4995}


\PY{n}{part} \PY{o}{=} \PY{p}{(}\PY{n}{N} \PY{o}{*} \PY{p}{(}\PY{n}{N} \PY{o}{\PYZhy{}} \PY{l+m+mi}{1}\PY{p}{)}\PY{p}{)} \PY{o}{/}\PY{o}{/} \PY{l+m+mi}{2} \PY{c+c1}{\PYZsh{}part = little \PYZsq{}n\PYZsq{} in nCr}
\PY{c+c1}{\PYZsh{}P\PYZus{}l = comb(part, L)*(p**L)*(1\PYZhy{}p)**(part\PYZhy{}L)}
\PY{n}{P\PYZus{}l} \PY{o}{=} \PY{n}{binom}\PY{o}{.}\PY{n}{pmf}\PY{p}{(}\PY{n}{L}\PY{p}{,} \PY{n}{part}\PY{p}{,} \PY{n}{p}\PY{p}{)}\PY{c+c1}{\PYZsh{}Finding P\PYZus{}l with the binom() function from scipy since comb() was throwing errors due to too large L}

\PY{c+c1}{\PYZsh{}Starting the plotting process}
\PY{n}{fig}\PY{p}{,} \PY{n}{ax} \PY{o}{=} \PY{n}{plt}\PY{o}{.}\PY{n}{subplots}\PY{p}{(}\PY{p}{)}\PY{c+c1}{\PYZsh{}Note: fig is the plot container, while ax (i.e., the \PYZsq{}axes\PYZsq{}) is the plot itself}
\PY{n}{ax}\PY{o}{.}\PY{n}{hist}\PY{p}{(}\PY{n}{num\PYZus{}edges\PYZus{}list}\PY{p}{,}\PY{n}{label}\PY{o}{=}\PY{l+s+s2}{\PYZdq{}}\PY{l+s+s2}{Networkx Generated Sample}\PY{l+s+s2}{\PYZdq{}}\PY{p}{,}\PY{n}{density}\PY{o}{=}\PY{k+kc}{True}\PY{p}{)}\PY{c+c1}{\PYZsh{}Plotting the networkx based random graph empirical data}
\PY{n}{ax}\PY{o}{.}\PY{n}{hist}\PY{p}{(}\PY{n}{my\PYZus{}num\PYZus{}edges\PYZus{}list}\PY{p}{,}\PY{n}{label}\PY{o}{=}\PY{l+s+s2}{\PYZdq{}}\PY{l+s+s2}{My Generated Sample}\PY{l+s+s2}{\PYZdq{}}\PY{p}{,}\PY{n}{density}\PY{o}{=}\PY{k+kc}{True}\PY{p}{)}\PY{c+c1}{\PYZsh{}Plotting my function based random graph empirical data}
\PY{n}{ax}\PY{o}{.}\PY{n}{plot}\PY{p}{(}\PY{n}{L}\PY{p}{,}\PY{n}{P\PYZus{}l}\PY{p}{,} \PY{n}{color} \PY{o}{=} \PY{l+s+s1}{\PYZsq{}}\PY{l+s+s1}{r}\PY{l+s+s1}{\PYZsq{}}\PY{p}{,}\PY{n}{label}\PY{o}{=}\PY{l+s+s2}{\PYZdq{}}\PY{l+s+s2}{Theoretical Distribution}\PY{l+s+s2}{\PYZdq{}}\PY{p}{)} \PY{c+c1}{\PYZsh{}Plotting the Theoretical}
\PY{c+c1}{\PYZsh{}Specifying figure details below:}
\PY{n}{handles}\PY{p}{,} \PY{n}{labels} \PY{o}{=} \PY{n}{ax}\PY{o}{.}\PY{n}{get\PYZus{}legend\PYZus{}handles\PYZus{}labels}\PY{p}{(}\PY{p}{)}
\PY{n}{ax}\PY{o}{.}\PY{n}{legend}\PY{p}{(}\PY{n}{handles}\PY{o}{=}\PY{n}{handles}\PY{p}{,} \PY{n}{labels}\PY{o}{=}\PY{n}{labels}\PY{p}{,}\PY{n}{loc}\PY{o}{=}\PY{l+s+s2}{\PYZdq{}}\PY{l+s+s2}{upper right}\PY{l+s+s2}{\PYZdq{}}\PY{p}{,}\PY{n}{fontsize}\PY{o}{=}\PY{l+s+s2}{\PYZdq{}}\PY{l+s+s2}{small}\PY{l+s+s2}{\PYZdq{}}\PY{p}{)}
\PY{n}{ax}\PY{o}{.}\PY{n}{set\PYZus{}title}\PY{p}{(}\PY{l+s+s2}{\PYZdq{}}\PY{l+s+s2}{Empirical vs. Theoretical Distribution of The Total Number of Edges: }\PY{l+s+s2}{\PYZsq{}}\PY{l+s+s2}{L}\PY{l+s+s2}{\PYZsq{}}\PY{l+s+s2}{\PYZdq{}}\PY{p}{)}
\PY{n}{ax}\PY{o}{.}\PY{n}{set\PYZus{}xlabel}\PY{p}{(}\PY{l+s+s2}{\PYZdq{}}\PY{l+s+s2}{L}\PY{l+s+s2}{\PYZdq{}}\PY{p}{)}
\PY{n}{ax}\PY{o}{.}\PY{n}{set\PYZus{}ylabel}\PY{p}{(}\PY{l+s+s2}{\PYZdq{}}\PY{l+s+s2}{Probability P\PYZus{}L}\PY{l+s+s2}{\PYZdq{}}\PY{p}{)}
\PY{n}{ax}\PY{o}{.}\PY{n}{set\PYZus{}xlim}\PY{p}{(}\PY{n}{np}\PY{o}{.}\PY{n}{min}\PY{p}{(}\PY{n}{L}\PY{p}{)}\PY{p}{,}\PY{n}{np}\PY{o}{.}\PY{n}{max}\PY{p}{(}\PY{n}{L}\PY{p}{)}\PY{p}{)}
\PY{n}{plt}\PY{o}{.}\PY{n}{show}\PY{p}{(}\PY{p}{)}
\end{Verbatim}
\end{tcolorbox}

    \begin{center}
    \adjustimage{max size={0.9\linewidth}{0.9\paperheight}}{Hughes_CNET5050_Assignment_1_Code_files/Hughes_CNET5050_Assignment_1_Code_3_0.png}
    \end{center}
    { \hspace*{\fill} \\}
    
    \paragraph{(ii) Generate NG = 100 independent graph realizations. For
every realization, record the degree of each node and concatenate the
results into a single sample containing 100 × n degree observations. Plot
the empirical degree distribution and overlay the
theoretical distribution derived in part (b). Clearly label the axes and
legend.}\label{ii-generate-ng-100-independent-graph-realizations.-for-every-realization-record-thedegree-of-each-node-and-concatenate-the-results-into-a-single-sample-containing-100-ndegree-observations.-plot-the-empirical-degree-distribution-and-overlay-the-theoreticaldistribution-derived-in-part-b.-clearly-label-the-axes-and-legend.}

    \begin{tcolorbox}[breakable, size=fbox, boxrule=1pt, pad at break*=1mm,colback=cellbackground, colframe=cellborder]
\prompt{In}{incolor}{18}{\boxspacing}
\begin{Verbatim}[commandchars=\\\{\}]
\PY{k+kn}{import}\PY{+w}{ }\PY{n+nn}{networkx}\PY{+w}{ }\PY{k}{as}\PY{+w}{ }\PY{n+nn}{nx}
\PY{k+kn}{import}\PY{+w}{ }\PY{n+nn}{matplotlib}\PY{n+nn}{.}\PY{n+nn}{pyplot}\PY{+w}{ }\PY{k}{as}\PY{+w}{ }\PY{n+nn}{plt}
\PY{k+kn}{from}\PY{+w}{ }\PY{n+nn}{scipy}\PY{n+nn}{.}\PY{n+nn}{stats}\PY{+w}{ }\PY{k+kn}{import} \PY{n}{binom}
\PY{k+kn}{from}\PY{+w}{ }\PY{n+nn}{scipy}\PY{n+nn}{.}\PY{n+nn}{special}\PY{+w}{ }\PY{k+kn}{import} \PY{n}{comb}

\PY{n}{N}\PY{o}{=}\PY{l+m+mi}{1000}
\PY{n}{p}\PY{o}{=}\PY{l+m+mf}{0.01}
\PY{n}{num\PYZus{}degree\PYZus{}list}\PY{o}{=}\PY{p}{[}\PY{p}{]}
\PY{n}{my\PYZus{}num\PYZus{}degree\PYZus{}list}\PY{o}{=}\PY{p}{[}\PY{p}{]}
\PY{k}{for} \PY{n}{i} \PY{o+ow}{in} \PY{n+nb}{range}\PY{p}{(}\PY{l+m+mi}{100}\PY{p}{)}\PY{p}{:}
  \PY{n}{G}\PY{o}{=}\PY{n}{nx}\PY{o}{.}\PY{n}{erdos\PYZus{}renyi\PYZus{}graph}\PY{p}{(}\PY{n}{N}\PY{p}{,} \PY{n}{p}\PY{p}{)}
  \PY{n}{degree} \PY{o}{=} \PY{n+nb}{dict}\PY{p}{(}\PY{n}{G}\PY{o}{.}\PY{n}{copy}\PY{p}{(}\PY{p}{)}\PY{o}{.}\PY{n}{degree}\PY{p}{(}\PY{p}{)}\PY{p}{)}
  \PY{n}{deg}\PY{o}{=} \PY{n+nb}{list}\PY{p}{(}\PY{n}{degree}\PY{o}{.}\PY{n}{values}\PY{p}{(}\PY{p}{)}\PY{p}{)}
  \PY{n}{num\PYZus{}degree\PYZus{}list}\PY{o}{.}\PY{n}{append}\PY{p}{(}\PY{n}{deg}\PY{p}{)}
  \PY{n}{G2}\PY{o}{=}\PY{n}{create\PYZus{}erdos\PYZus{}renyi\PYZus{}random\PYZus{}graph}\PY{p}{(}\PY{n}{N}\PY{p}{,} \PY{n}{p}\PY{p}{)}
  \PY{n}{degree} \PY{o}{=} \PY{n+nb}{dict}\PY{p}{(}\PY{n}{G2}\PY{o}{.}\PY{n}{copy}\PY{p}{(}\PY{p}{)}\PY{o}{.}\PY{n}{degree}\PY{p}{(}\PY{p}{)}\PY{p}{)}
  \PY{n}{deg}\PY{o}{=} \PY{n+nb}{list}\PY{p}{(}\PY{n}{degree}\PY{o}{.}\PY{n}{values}\PY{p}{(}\PY{p}{)}\PY{p}{)}
  \PY{n}{my\PYZus{}num\PYZus{}degree\PYZus{}list}\PY{o}{.}\PY{n}{append}\PY{p}{(}\PY{n}{deg}\PY{p}{)}

\PY{n}{k} \PY{o}{=} \PY{n}{np}\PY{o}{.}\PY{n}{arange}\PY{p}{(}\PY{l+m+mf}{9.99}\PY{o}{\PYZhy{}}\PY{l+m+mf}{9.99}\PY{p}{,}\PY{l+m+mf}{9.99}\PY{o}{+}\PY{l+m+mf}{9.99}\PY{o}{+}\PY{l+m+mf}{0.01}\PY{p}{,}\PY{n}{step}\PY{o}{=}\PY{l+m+mf}{0.01}\PY{p}{,}\PY{n}{dtype} \PY{o}{=} \PY{n+nb}{float}\PY{p}{)}\PY{c+c1}{\PYZsh{}Picking a handful of values centered around the average I calculated earlier of: 9.99}

\PY{n}{part} \PY{o}{=} \PY{p}{(}\PY{n}{N} \PY{o}{\PYZhy{}} \PY{l+m+mi}{1}\PY{p}{)} \PY{c+c1}{\PYZsh{}part = little \PYZsq{}n\PYZsq{} in nCr}
\PY{n}{P\PYZus{}l} \PY{o}{=} \PY{n}{comb}\PY{p}{(}\PY{n}{part}\PY{p}{,} \PY{n}{k}\PY{p}{)}\PY{o}{*}\PY{p}{(}\PY{n}{p}\PY{o}{*}\PY{o}{*}\PY{n}{k}\PY{p}{)}\PY{o}{*}\PY{p}{(}\PY{l+m+mi}{1}\PY{o}{\PYZhy{}}\PY{n}{p}\PY{p}{)}\PY{o}{*}\PY{o}{*}\PY{p}{(}\PY{n}{part}\PY{o}{\PYZhy{}}\PY{n}{k}\PY{p}{)}\PY{c+c1}{\PYZsh{}Calculating P\PYZus{}l with comb()}
\PY{c+c1}{\PYZsh{}P\PYZus{}l = binom.pmf(k, part, p) \PYZsh{}Unfortunately because nCr requires integers as input, I can only calculate P\PYZus{}l for integer numbers with binom()}
\PY{n+nb}{print}\PY{p}{(}\PY{l+s+s2}{\PYZdq{}}\PY{l+s+s2}{P\PYZus{}l: }\PY{l+s+s2}{\PYZdq{}}\PY{p}{,}\PY{n}{P\PYZus{}l}\PY{p}{)}
\PY{n+nb}{print}\PY{p}{(}\PY{l+s+s2}{\PYZdq{}}\PY{l+s+s2}{k: }\PY{l+s+s2}{\PYZdq{}}\PY{p}{,}\PY{n}{k}\PY{p}{)}

\PY{c+c1}{\PYZsh{}Starting the plotting process}
\PY{n}{fig}\PY{p}{,} \PY{n}{ax} \PY{o}{=} \PY{n}{plt}\PY{o}{.}\PY{n}{subplots}\PY{p}{(}\PY{p}{)} \PY{c+c1}{\PYZsh{}Note: fig is the plot container, while ax (i.e., the \PYZsq{}axes\PYZsq{}) is the plot itself}
\PY{n}{ax}\PY{o}{.}\PY{n}{hist}\PY{p}{(}\PY{n}{num\PYZus{}degree\PYZus{}list}\PY{p}{,}\PY{n}{label}\PY{o}{=}\PY{l+s+s2}{\PYZdq{}}\PY{l+s+s2}{Networkx Generated Sample}\PY{l+s+s2}{\PYZdq{}}\PY{p}{,}\PY{n}{density}\PY{o}{=}\PY{k+kc}{True}\PY{p}{,}\PY{n}{color} \PY{o}{=} \PY{l+s+s1}{\PYZsq{}}\PY{l+s+s1}{b}\PY{l+s+s1}{\PYZsq{}}\PY{p}{,}\PY{n}{histtype}\PY{o}{=}\PY{l+s+s1}{\PYZsq{}}\PY{l+s+s1}{stepfilled}\PY{l+s+s1}{\PYZsq{}}\PY{p}{)}\PY{c+c1}{\PYZsh{}Plotting Networkx based empirical}
\PY{n}{ax}\PY{o}{.}\PY{n}{hist}\PY{p}{(}\PY{n}{my\PYZus{}num\PYZus{}degree\PYZus{}list}\PY{p}{,}\PY{n}{label}\PY{o}{=}\PY{l+s+s2}{\PYZdq{}}\PY{l+s+s2}{My Generated Sample}\PY{l+s+s2}{\PYZdq{}}\PY{p}{,}\PY{n}{density}\PY{o}{=}\PY{k+kc}{True}\PY{p}{,}\PY{n}{color} \PY{o}{=} \PY{l+s+s1}{\PYZsq{}}\PY{l+s+s1}{orange}\PY{l+s+s1}{\PYZsq{}}\PY{p}{,}\PY{n}{histtype}\PY{o}{=}\PY{l+s+s1}{\PYZsq{}}\PY{l+s+s1}{stepfilled}\PY{l+s+s1}{\PYZsq{}}\PY{p}{)}\PY{c+c1}{\PYZsh{}Plotting My Random Graph based empirical}
\PY{n}{ax}\PY{o}{.}\PY{n}{plot}\PY{p}{(}\PY{n}{k}\PY{p}{,}\PY{n}{P\PYZus{}l}\PY{p}{,} \PY{n}{color} \PY{o}{=} \PY{l+s+s1}{\PYZsq{}}\PY{l+s+s1}{r}\PY{l+s+s1}{\PYZsq{}}\PY{p}{,}\PY{n}{label}\PY{o}{=}\PY{l+s+s2}{\PYZdq{}}\PY{l+s+s2}{Theoretical Distribution}\PY{l+s+s2}{\PYZdq{}}\PY{p}{)}   \PY{c+c1}{\PYZsh{}Plotting the Theoretical}

\PY{c+c1}{\PYZsh{}Specifying figure details below:}
\PY{n}{handles}\PY{p}{,} \PY{n}{labels} \PY{o}{=} \PY{n}{ax}\PY{o}{.}\PY{n}{get\PYZus{}legend\PYZus{}handles\PYZus{}labels}\PY{p}{(}\PY{p}{)}
\PY{n}{ax}\PY{o}{.}\PY{n}{legend}\PY{p}{(}\PY{n}{handles}\PY{o}{=}\PY{n}{handles}\PY{p}{,} \PY{n}{labels}\PY{o}{=}\PY{n}{labels}\PY{p}{,}\PY{n}{loc}\PY{o}{=}\PY{l+s+s2}{\PYZdq{}}\PY{l+s+s2}{upper right}\PY{l+s+s2}{\PYZdq{}}\PY{p}{,}\PY{n}{fontsize}\PY{o}{=}\PY{l+s+s2}{\PYZdq{}}\PY{l+s+s2}{small}\PY{l+s+s2}{\PYZdq{}}\PY{p}{)}
\PY{n}{ax}\PY{o}{.}\PY{n}{set\PYZus{}title}\PY{p}{(}\PY{l+s+s2}{\PYZdq{}}\PY{l+s+s2}{Empirical vs. Theoretical Distribution of The Average Degree: }\PY{l+s+s2}{\PYZsq{}}\PY{l+s+s2}{k}\PY{l+s+s2}{\PYZsq{}}\PY{l+s+s2}{\PYZdq{}}\PY{p}{)}
\PY{n}{ax}\PY{o}{.}\PY{n}{set\PYZus{}xlabel}\PY{p}{(}\PY{l+s+s2}{\PYZdq{}}\PY{l+s+s2}{k}\PY{l+s+s2}{\PYZdq{}}\PY{p}{)}
\PY{n}{ax}\PY{o}{.}\PY{n}{set\PYZus{}ylabel}\PY{p}{(}\PY{l+s+s2}{\PYZdq{}}\PY{l+s+s2}{Probability P\PYZus{}k}\PY{l+s+s2}{\PYZdq{}}\PY{p}{)}
\PY{n}{ax}\PY{o}{.}\PY{n}{set\PYZus{}xlim}\PY{p}{(}\PY{n}{np}\PY{o}{.}\PY{n}{min}\PY{p}{(}\PY{n}{k}\PY{p}{)}\PY{p}{,}\PY{n}{np}\PY{o}{.}\PY{n}{max}\PY{p}{(}\PY{n}{k}\PY{p}{)}\PY{p}{)}
\PY{n}{plt}\PY{o}{.}\PY{n}{show}\PY{p}{(}\PY{p}{)}
\end{Verbatim}
\end{tcolorbox}

    \begin{Verbatim}[commandchars=\\\{\}]
P\_l:  [4.36073206e-05 4.48819923e-05 4.61864302e-05 {\ldots} 1.81293599e-03
 1.79979733e-03 1.78674498e-03]
k:  [0.000e+00 1.000e-02 2.000e-02 {\ldots} 1.997e+01 1.998e+01 1.999e+01]
    \end{Verbatim}

    \begin{center}
    \adjustimage{max size={0.9\linewidth}{0.9\paperheight}}{Hughes_CNET5050_Assignment_1_Code_files/Hughes_CNET5050_Assignment_1_Code_5_1.png}
    \end{center}
    { \hspace*{\fill} \\}
    
    \paragraph{II. Real-world networks {[}25
marks{]}\newline}\label{ii.-real-world-networks-25-marks}

Download the airport-network.gml dataset from
https://tinyurl.com/4rfrx4p9. Using the NetworkX library in Python, load
the dataset with nx.read\_gml and complete the following analyses.

\paragraph{Report the number of nodes n, the number of edges L, the
average degree ⟨k⟩, and the average local clustering coefficient of the
network.}\label{report-the-number-of-nodes-n-the-number-of-edges-l-the-average-degree-k-and-the-average-local-clustering-coefficient-of-the-network.}

    \begin{tcolorbox}[breakable, size=fbox, boxrule=1pt, pad at break*=1mm,colback=cellbackground, colframe=cellborder]
\prompt{In}{incolor}{19}{\boxspacing}
\begin{Verbatim}[commandchars=\\\{\}]
\PY{c+c1}{\PYZsh{}Reading in the Data}
\PY{n}{Airport} \PY{o}{=} \PY{n}{nx}\PY{o}{.}\PY{n}{read\PYZus{}gml}\PY{p}{(}\PY{l+s+s1}{\PYZsq{}}\PY{l+s+s1}{C:/Users/rh614/OneDrive/Documents/ALL REPOS/Python/Fundamentals\PYZhy{}of\PYZhy{}Complex\PYZhy{}Networks/Assignment\PYZus{}1/data/airport\PYZhy{}network.gml}\PY{l+s+s1}{\PYZsq{}}\PY{p}{)}
\PY{n}{N} \PY{o}{=} \PY{n}{Airport}\PY{o}{.}\PY{n}{number\PYZus{}of\PYZus{}nodes}\PY{p}{(}\PY{p}{)}
\PY{n+nb}{print}\PY{p}{(}\PY{l+s+s2}{\PYZdq{}}\PY{l+s+s2}{N: }\PY{l+s+s2}{\PYZdq{}}\PY{p}{,}\PY{n}{N}\PY{p}{)}
\PY{n}{L} \PY{o}{=} \PY{n}{Airport}\PY{o}{.}\PY{n}{number\PYZus{}of\PYZus{}edges}\PY{p}{(}\PY{p}{)}
\PY{n+nb}{print}\PY{p}{(}\PY{l+s+s2}{\PYZdq{}}\PY{l+s+s2}{L: }\PY{l+s+s2}{\PYZdq{}}\PY{p}{,}\PY{n}{L}\PY{p}{)}
\PY{n}{avg\PYZus{}deg}\PY{o}{=} \PY{n}{np}\PY{o}{.}\PY{n}{mean}\PY{p}{(}\PY{n+nb}{list}\PY{p}{(}\PY{n+nb}{dict}\PY{p}{(}\PY{n}{Airport}\PY{o}{.}\PY{n}{degree}\PY{p}{(}\PY{p}{)}\PY{p}{)}\PY{o}{.}\PY{n}{values}\PY{p}{(}\PY{p}{)}\PY{p}{)}\PY{p}{)}
\PY{n+nb}{print}\PY{p}{(}\PY{l+s+s2}{\PYZdq{}}\PY{l+s+s2}{Average Degree: }\PY{l+s+s2}{\PYZdq{}}\PY{p}{,}\PY{n}{avg\PYZus{}deg}\PY{p}{)}
\PY{n}{C\PYZus{}i} \PY{o}{=} \PY{n}{nx}\PY{o}{.}\PY{n}{average\PYZus{}clustering}\PY{p}{(}\PY{n}{Airport}\PY{p}{)}
\PY{n+nb}{print}\PY{p}{(}\PY{l+s+s2}{\PYZdq{}}\PY{l+s+s2}{C\PYZus{}i: }\PY{l+s+s2}{\PYZdq{}}\PY{p}{,}\PY{n}{C\PYZus{}i}\PY{p}{)}
\end{Verbatim}
\end{tcolorbox}

    \begin{Verbatim}[commandchars=\\\{\}]
N:  1574
L:  17215
Average Degree:  21.874205844980942
C\_i:  0.504152539209577
    \end{Verbatim}

    \paragraph{Plot the empirical degree distribution using logarithmic
scales on both axes. Bin the node degrees and plot a normalized
histogram. Clearly label the axes and briefly comment on the shape of
the
distribution.}\label{plot-the-empirical-degree-distribution-using-logarithmic-scales-on-both-axes.-bin-the-node-degrees-and-plot-a-normalized-histogram.-clearly-label-the-axes-and-briefly-comment-on-the-shape-of-the-distribution.}

    \begin{tcolorbox}[breakable, size=fbox, boxrule=1pt, pad at break*=1mm,colback=cellbackground, colframe=cellborder]
\prompt{In}{incolor}{20}{\boxspacing}
\begin{Verbatim}[commandchars=\\\{\}]
\PY{n}{deg} \PY{o}{=} \PY{n+nb}{list}\PY{p}{(}\PY{n+nb}{dict}\PY{p}{(}\PY{n}{Airport}\PY{o}{.}\PY{n}{degree}\PY{p}{(}\PY{p}{)}\PY{p}{)}\PY{o}{.}\PY{n}{values}\PY{p}{(}\PY{p}{)}\PY{p}{)} \PY{c+c1}{\PYZsh{}Gathering the degree values as a list}
\PY{c+c1}{\PYZsh{}Printing the Min and Max Degree for reference purposes}
\PY{n+nb}{print}\PY{p}{(}\PY{l+s+s2}{\PYZdq{}}\PY{l+s+s2}{Max Deg: }\PY{l+s+s2}{\PYZdq{}}\PY{p}{,}\PY{n}{np}\PY{o}{.}\PY{n}{max}\PY{p}{(}\PY{n}{deg}\PY{p}{)}\PY{p}{)}
\PY{n+nb}{print}\PY{p}{(}\PY{l+s+s2}{\PYZdq{}}\PY{l+s+s2}{Min Deg: }\PY{l+s+s2}{\PYZdq{}}\PY{p}{,}\PY{n}{np}\PY{o}{.}\PY{n}{min}\PY{p}{(}\PY{n}{deg}\PY{p}{)}\PY{p}{)}

\PY{c+c1}{\PYZsh{}Starting the plotting process}
\PY{n}{fig}\PY{p}{,} \PY{n}{ax} \PY{o}{=} \PY{n}{plt}\PY{o}{.}\PY{n}{subplots}\PY{p}{(}\PY{p}{)}
\PY{n}{bins} \PY{o}{=} \PY{n}{np}\PY{o}{.}\PY{n}{geomspace}\PY{p}{(}\PY{n}{np}\PY{o}{.}\PY{n}{min}\PY{p}{(}\PY{n}{deg}\PY{p}{)}\PY{p}{,}\PY{n}{np}\PY{o}{.}\PY{n}{max}\PY{p}{(}\PY{n}{deg}\PY{p}{)}\PY{p}{)} \PY{c+c1}{\PYZsh{}Making bins that are log\PYZhy{}scaled}

\PY{c+c1}{\PYZsh{}Plotting the degree distribution histogram using log\PYZhy{}scaled bins}
\PY{n}{ax}\PY{o}{.}\PY{n}{hist}\PY{p}{(}\PY{n}{deg}\PY{p}{,}\PY{n}{bins}\PY{o}{=}\PY{n}{bins}\PY{p}{,}\PY{n}{label}\PY{o}{=}\PY{l+s+s2}{\PYZdq{}}\PY{l+s+s2}{Airport Sample}\PY{l+s+s2}{\PYZdq{}}\PY{p}{,}\PY{n}{density}\PY{o}{=}\PY{k+kc}{True}\PY{p}{,}\PY{n}{color} \PY{o}{=} \PY{l+s+s1}{\PYZsq{}}\PY{l+s+s1}{b}\PY{l+s+s1}{\PYZsq{}}\PY{p}{,}\PY{n}{histtype}\PY{o}{=}\PY{l+s+s1}{\PYZsq{}}\PY{l+s+s1}{step}\PY{l+s+s1}{\PYZsq{}}\PY{p}{,}\PY{n}{log}\PY{o}{=}\PY{k+kc}{True}\PY{p}{)}
\PY{n}{ax}\PY{o}{.}\PY{n}{loglog}\PY{p}{(}\PY{p}{)}\PY{c+c1}{\PYZsh{}Making the axes log\PYZhy{}log}

\PY{c+c1}{\PYZsh{}Specifying figure details below:}
\PY{n}{handles}\PY{p}{,} \PY{n}{labels} \PY{o}{=} \PY{n}{ax}\PY{o}{.}\PY{n}{get\PYZus{}legend\PYZus{}handles\PYZus{}labels}\PY{p}{(}\PY{p}{)}
\PY{n}{ax}\PY{o}{.}\PY{n}{legend}\PY{p}{(}\PY{n}{handles}\PY{o}{=}\PY{n}{handles}\PY{p}{,} \PY{n}{labels}\PY{o}{=}\PY{n}{labels}\PY{p}{,}\PY{n}{loc}\PY{o}{=}\PY{l+s+s2}{\PYZdq{}}\PY{l+s+s2}{upper right}\PY{l+s+s2}{\PYZdq{}}\PY{p}{,}\PY{n}{fontsize}\PY{o}{=}\PY{l+s+s2}{\PYZdq{}}\PY{l+s+s2}{small}\PY{l+s+s2}{\PYZdq{}}\PY{p}{)}
\PY{n}{ax}\PY{o}{.}\PY{n}{set\PYZus{}title}\PY{p}{(}\PY{l+s+s2}{\PYZdq{}}\PY{l+s+s2}{Empirical Degree Distribution of the Airport Network}\PY{l+s+s2}{\PYZdq{}}\PY{p}{)}
\PY{n}{ax}\PY{o}{.}\PY{n}{set\PYZus{}xlabel}\PY{p}{(}\PY{l+s+s2}{\PYZdq{}}\PY{l+s+s2}{L}\PY{l+s+s2}{\PYZdq{}}\PY{p}{)}
\PY{n}{ax}\PY{o}{.}\PY{n}{set\PYZus{}ylabel}\PY{p}{(}\PY{l+s+s2}{\PYZdq{}}\PY{l+s+s2}{Probability P\PYZus{}L}\PY{l+s+s2}{\PYZdq{}}\PY{p}{)}
\PY{n}{plt}\PY{o}{.}\PY{n}{show}\PY{p}{(}\PY{p}{)}
\end{Verbatim}
\end{tcolorbox}

    \begin{Verbatim}[commandchars=\\\{\}]
Max Deg:  314
Min Deg:  1
    \end{Verbatim}

    \begin{center}
    \adjustimage{max size={0.9\linewidth}{0.9\paperheight}}{Hughes_CNET5050_Assignment_1_Code_files/Hughes_CNET5050_Assignment_1_Code_9_1.png}
    \end{center}
    { \hspace*{\fill} \\}
    
    The shape of the distribution is approximately a straight line with
negative slope on the log-log scale. This indicates the data follows a power-law or truncated power-law distribution. The range of the degree distribution varies from
1 to 314, with nodes having a higher probability of having a degree that is lower (e.g, less
than 10). 1 being the most probable degree to occur for a given node, and 314 the least probable to occur.

    \paragraph{Overlay a Poisson degree distribution with the same mean
degree ⟨k⟩ on the empirical degree distribution. Compare the two
distributions and briefly discuss their main differences.}\label{overlay-a-poisson-degree-distribution-with-the-same-mean-degree-k-on-the-empirical-degree-distribution.-compare-the-two-distributions-and-briefly-discuss-their-main-differences.part-2-50-marks}

    \begin{tcolorbox}[breakable, size=fbox, boxrule=1pt, pad at break*=1mm,colback=cellbackground, colframe=cellborder]
\prompt{In}{incolor}{21}{\boxspacing}
\begin{Verbatim}[commandchars=\\\{\}]
\PY{k+kn}{from}\PY{+w}{ }\PY{n+nn}{scipy}\PY{n+nn}{.}\PY{n+nn}{stats}\PY{+w}{ }\PY{k+kn}{import} \PY{n}{poisson}
\PY{n}{deg} \PY{o}{=} \PY{n+nb}{list}\PY{p}{(}\PY{n+nb}{dict}\PY{p}{(}\PY{n}{Airport}\PY{o}{.}\PY{n}{degree}\PY{p}{(}\PY{p}{)}\PY{p}{)}\PY{o}{.}\PY{n}{values}\PY{p}{(}\PY{p}{)}\PY{p}{)}
\PY{c+c1}{\PYZsh{}Printing the Min and Max Degree for reference purposes}
\PY{n+nb}{print}\PY{p}{(}\PY{l+s+s2}{\PYZdq{}}\PY{l+s+s2}{Max Deg: }\PY{l+s+s2}{\PYZdq{}}\PY{p}{,}\PY{n}{np}\PY{o}{.}\PY{n}{max}\PY{p}{(}\PY{n}{deg}\PY{p}{)}\PY{p}{)}
\PY{n+nb}{print}\PY{p}{(}\PY{l+s+s2}{\PYZdq{}}\PY{l+s+s2}{Min Deg: }\PY{l+s+s2}{\PYZdq{}}\PY{p}{,}\PY{n}{np}\PY{o}{.}\PY{n}{min}\PY{p}{(}\PY{n}{deg}\PY{p}{)}\PY{p}{)}
\PY{n}{avg\PYZus{}deg}\PY{o}{=} \PY{n}{np}\PY{o}{.}\PY{n}{mean}\PY{p}{(}\PY{n+nb}{list}\PY{p}{(}\PY{n+nb}{dict}\PY{p}{(}\PY{n}{Airport}\PY{o}{.}\PY{n}{degree}\PY{p}{(}\PY{p}{)}\PY{p}{)}\PY{o}{.}\PY{n}{values}\PY{p}{(}\PY{p}{)}\PY{p}{)}\PY{p}{)}\PY{c+c1}{\PYZsh{}Getting the average degree for the Poisson}

\PY{c+c1}{\PYZsh{}Starting the plotting process}
\PY{n}{fig}\PY{p}{,} \PY{n}{ax} \PY{o}{=} \PY{n}{plt}\PY{o}{.}\PY{n}{subplots}\PY{p}{(}\PY{p}{)}
\PY{n}{bins} \PY{o}{=} \PY{n}{np}\PY{o}{.}\PY{n}{geomspace}\PY{p}{(}\PY{n}{np}\PY{o}{.}\PY{n}{min}\PY{p}{(}\PY{n}{deg}\PY{p}{)}\PY{p}{,}\PY{n}{np}\PY{o}{.}\PY{n}{max}\PY{p}{(}\PY{n}{deg}\PY{p}{)}\PY{p}{,}\PY{n}{num}\PY{o}{=}\PY{l+m+mi}{200}\PY{p}{)}\PY{c+c1}{\PYZsh{}Making log scaled bins for the empirical data}

\PY{c+c1}{\PYZsh{}Note mu=lambda, k=x}
\PY{n}{k\PYZus{}values} \PY{o}{=} \PY{n+nb}{range}\PY{p}{(}\PY{n}{np}\PY{o}{.}\PY{n}{min}\PY{p}{(}\PY{n}{deg}\PY{p}{)}\PY{p}{,}\PY{n}{np}\PY{o}{.}\PY{n}{max}\PY{p}{(}\PY{n}{deg}\PY{p}{)}\PY{o}{+}\PY{l+m+mi}{1}\PY{p}{,}\PY{l+m+mi}{1}\PY{p}{)}\PY{c+c1}{\PYZsh{}Generating the input to the poisson PDF as all integers of k from empirical\PYZhy{}data}
\PY{n}{poisson} \PY{o}{=} \PY{n}{poisson}\PY{o}{.}\PY{n}{pmf}\PY{p}{(}\PY{n}{k}\PY{o}{=}\PY{n}{k\PYZus{}values}\PY{p}{,}\PY{n}{mu}\PY{o}{=}\PY{n}{avg\PYZus{}deg}\PY{p}{)}\PY{c+c1}{\PYZsh{}Getting the P\PYZus{}k based on Poisson}

\PY{c+c1}{\PYZsh{}Plotting the Empirical and Poisson Distributions}
\PY{n}{ax}\PY{o}{.}\PY{n}{hist}\PY{p}{(}\PY{n}{deg}\PY{p}{,}\PY{n}{bins}\PY{o}{=}\PY{n}{bins}\PY{p}{,}\PY{n}{label}\PY{o}{=}\PY{l+s+s2}{\PYZdq{}}\PY{l+s+s2}{Airport Sample}\PY{l+s+s2}{\PYZdq{}}\PY{p}{,}\PY{n}{density}\PY{o}{=}\PY{k+kc}{True}\PY{p}{,}\PY{n}{color} \PY{o}{=} \PY{l+s+s1}{\PYZsq{}}\PY{l+s+s1}{b}\PY{l+s+s1}{\PYZsq{}}\PY{p}{,}\PY{n}{histtype}\PY{o}{=}\PY{l+s+s1}{\PYZsq{}}\PY{l+s+s1}{step}\PY{l+s+s1}{\PYZsq{}}\PY{p}{,}\PY{n}{log}\PY{o}{=}\PY{k+kc}{True}\PY{p}{)}\PY{c+c1}{\PYZsh{}Plotting the Empirical}
\PY{n}{ax}\PY{o}{.}\PY{n}{loglog}\PY{p}{(}\PY{p}{)}\PY{c+c1}{\PYZsh{}Setting to log\PYZhy{}log scale}
\PY{n}{ax}\PY{o}{.}\PY{n}{plot}\PY{p}{(}\PY{n}{k\PYZus{}values}\PY{p}{,}\PY{n}{poisson}\PY{p}{,}\PY{n}{label}\PY{o}{=}\PY{l+s+s2}{\PYZdq{}}\PY{l+s+s2}{Poisson Distribution}\PY{l+s+s2}{\PYZdq{}}\PY{p}{,}\PY{n}{color} \PY{o}{=} \PY{l+s+s1}{\PYZsq{}}\PY{l+s+s1}{r}\PY{l+s+s1}{\PYZsq{}}\PY{p}{)}\PY{c+c1}{\PYZsh{}Plotting the Poisson}

\PY{c+c1}{\PYZsh{}Specifying figure details below:}
\PY{n}{handles}\PY{p}{,} \PY{n}{labels} \PY{o}{=} \PY{n}{ax}\PY{o}{.}\PY{n}{get\PYZus{}legend\PYZus{}handles\PYZus{}labels}\PY{p}{(}\PY{p}{)}\PY{c+c1}{\PYZsh{}Gathering the labels to make legends per plot}
\PY{n}{ax}\PY{o}{.}\PY{n}{legend}\PY{p}{(}\PY{n}{handles}\PY{o}{=}\PY{n}{handles}\PY{p}{,} \PY{n}{labels}\PY{o}{=}\PY{n}{labels}\PY{p}{,}\PY{n}{loc}\PY{o}{=}\PY{l+s+s2}{\PYZdq{}}\PY{l+s+s2}{upper right}\PY{l+s+s2}{\PYZdq{}}\PY{p}{,}\PY{n}{fontsize}\PY{o}{=}\PY{l+s+s2}{\PYZdq{}}\PY{l+s+s2}{small}\PY{l+s+s2}{\PYZdq{}}\PY{p}{)}\PY{c+c1}{\PYZsh{}Adding legend}
\PY{n}{ax}\PY{o}{.}\PY{n}{set\PYZus{}title}\PY{p}{(}\PY{l+s+s2}{\PYZdq{}}\PY{l+s+s2}{Empirical Degree Distribution of the Airport Network vs Poisson}\PY{l+s+s2}{\PYZdq{}}\PY{p}{)}
\PY{n}{ax}\PY{o}{.}\PY{n}{set\PYZus{}xlabel}\PY{p}{(}\PY{l+s+s2}{\PYZdq{}}\PY{l+s+s2}{k}\PY{l+s+s2}{\PYZdq{}}\PY{p}{)}
\PY{n}{ax}\PY{o}{.}\PY{n}{set\PYZus{}ylim}\PY{p}{(}\PY{l+m+mf}{10e\PYZhy{}26}\PY{p}{,}\PY{l+m+mi}{2}\PY{p}{)}
\PY{n}{ax}\PY{o}{.}\PY{n}{set\PYZus{}ylabel}\PY{p}{(}\PY{l+s+s2}{\PYZdq{}}\PY{l+s+s2}{Probability P\PYZus{}k}\PY{l+s+s2}{\PYZdq{}}\PY{p}{)}
\PY{n}{plt}\PY{o}{.}\PY{n}{show}\PY{p}{(}\PY{p}{)}
\end{Verbatim}
\end{tcolorbox}

    \begin{Verbatim}[commandchars=\\\{\}]
Max Deg:  314
Min Deg:  1
    \end{Verbatim}

    \begin{center}
    \adjustimage{max size={0.9\linewidth}{0.9\paperheight}}{Hughes_CNET5050_Assignment_1_Code_files/Hughes_CNET5050_Assignment_1_Code_12_1.png}
    \end{center}
    { \hspace*{\fill} \\}
    
    Notably, the main difference is that the airport samples distribution
follows closer to a straight line (i.e., a power law distribution, or
truncated power law distribution) while the Poisson Distribution on the
log-log scale doesn't represent the data well. The Poisson sharply cuts
down in probability before the degree reaches 100, while the airport
clearly still has many nodes (\emph{relatively speaking}) with degree
higher than 100.

Acknowledgements: \newline
AI was not used in the creation of this assignment. I utilized stack-overflow for guidance, especially on how to generate
the Binomial distribution and Poisson distribution respectively. I also referred to our coding lectures in CNET5051 and
Python documentation pages for guidance on specific packages (e.g., Numpy, NetworkX, and SciPy) to understand how to leverage
relevant commands correctly, resolve bugs and fix errors in my code.
    
    
\end{document}
